\section*{Chapitre \ref{ch:m1:asian-and-emerging-equity-markets} --- Marchés d'actions asiatiques et émergents}

\begin{solution}{exo:m1:asian-and-emerging-equity-markets:1}
Limite haute 26.40, limite basse 21.60. Les deux coïncident ici parce que
$24.00\times1.1$ tombe sur un tick. En général $P\times1.1$ n'y tombe pas,
et les bourses diffèrent dans l'arrondi au tick le plus proche, vers
l'intérieur ou selon une règle qui leur est propre~: un système qui arrondit
autrement envoie des ordres à un tick hors de la bande, qui sont rejetés précisément les jours où cela compte.
\end{solution}

\begin{solution}{exo:m1:asian-and-emerging-equity-markets:2}
$250\,000/61.30 = 4\,078$ actions, soit 8 lots complets~: 4\,000 actions, HKD
245\,200~; droit $0.1\% = $ HKD 245.20 (et autant à la vente finale). HKD
4\,800 restent non investis~: la taille de lot compte pour les petits comptes et
pour les couvertures précises.
\end{solution}

\begin{solution}{exo:m1:asian-and-emerging-equity-markets:3}
(a) $0.1\%$ de chaque côté~: INR 5\,000 $+$ 5\,000 $=$ 10\,000. (b) $0.05\%
\times 20$ millions de roupies $=$ INR 10\,000, à la vente seulement. (c) $0.15\% \times
400\,000 =$ INR 600.
\end{solution}

\begin{solution}{exo:m1:asian-and-emerging-equity-markets:4}
$\ln0.65/\ln0.9 = 4.09$~: cinq jours, quatre bloqués. À 20~\%~: $1.93$,
deux jours, un bloqué. À 5~\%~: $8.40$, neuf jours, huit bloqués.
\end{solution}

\begin{solution}{exo:m1:asian-and-emerging-equity-markets:5}
Sans taxe~: $h^\star = 6/4 = 1.5$ jour. Avec le droit~: $26/4 = 6.5$ jours. À dix
jours, l'avantage brut est de 40 points de base~: 34 nets sans le droit, 14 avec lui
--- la taxe prélève 59~\% du profit d'une stratégie qu'elle n'interdit pas.
\end{solution}

\begin{solution}{exo:m1:asian-and-emerging-equity-markets:6}
Elle ne peut pas vendre les actions. Elle peut vendre des contrats à terme sur
l'indice, ou sur l'action si un contrat à terme sur action existe, pour le même notionnel~: cela élimine
la composante de marché du risque et laisse le mouvement propre de l'action jusqu'à
l'ouverture du lendemain, plus la base de la couverture. Si elle détenait déjà la même
action depuis un jour précédent, elle pourrait vendre \emph{ces} actions aujourd'hui~: l'interdiction
porte sur les achats de la journée, c'est pourquoi les fonds locaux conservent un stock permanent
dans les valeurs qu'ils négocient en cours de séance.
\end{solution}

\begin{solution}{exo:m1:asian-and-emerging-equity-markets:7}
Écart type et autocorrélation~: 2.60~\% et $+0.195$ à 5~\%~; 3.03~\%
et $+0.088$ à 10~\%~; 3.26~\% et $+0.033$ à 20~\%, contre une vraie valeur de 3.45~\% et
zéro. Plus la limite est serrée, plus une part de la distribution est tronquée et reportée
au jour suivant~: la volatilité est sous-estimée et un faux momentum croît à peu près
en proportion de la part des jours qui finissent à une limite.
\end{solution}

\begin{solution}{exo:m1:asian-and-emerging-equity-markets:8}
D'abord, la clôture d'un jour à la hausse limite n'est pas un prix auquel on
peut acheter~: le carnet d'ordres montre une file d'acheteurs et presque aucun
vendeur, et l'exécution espérée d'un nouvel ordre est proche de zéro --- et
elle est maximale précisément dans les cas où le lendemain est mauvais
(lorsque les détenteurs acceptaient de vendre à la limite). Ensuite, l'achat
ne pouvait pas être revendu le jour même, si bien que la sortie «~du lendemain~»
peut elle-même être bloquée si l'action ouvre en baisse. La statistique à publier
est le rendement conditionnel à un ordre \emph{effectivement exécuté} à la limite,
pondéré par la quantité exécutée~: une moyenne pondérée par les exécutions, et non une moyenne à pondération égale.
\end{solution}

\begin{solution}{pb:m1:asian-and-emerging-equity-markets:1}
\textbf{1.} $29.00/(1 + 0.04\times0.25) = 28.71$.
\textbf{2.} 22.00~; 24.20~; 26.62~; 29.28.
\textbf{3.} Vendredi~: 29.28 est la première limite au-dessus de 28.71. Bloqué le mardi,
le mercredi et le jeudi~: trois jours.
\textbf{4.} $\ln(1.4356)/\ln(1.1) = 3.79$~; $\lceil3.79\rceil - 1 = 3$.
\textbf{5.} Aucune~: 9 millions d'actions sont devant vous et 300\,000 sont échangées.
\textbf{6.} $28.71/22.00 - 1 = 30.5\%$.
\textbf{7.} $100\,000 \times (28.71 - 22.00) = 671\,000$.
\textbf{8.} Priorité temporelle à l'ouverture~: les ordres saisis dans les
premières microsecondes après l'ouverture de la saisie, ou mis en file pendant la pré-ouverture par les
courtiers disposant des liaisons les plus rapides. Une société paierait jusqu'au gain espéré
multiplié par sa probabilité d'être exécutée --- ce qui fait des files à la hausse limite
une course à la latence sur les marchés soumis à des limites de cours.
\textbf{9.} Le mercredi au plus tôt, à la limite de mercredi de 24.20 au mieux
--- et seulement si quelqu'un achète à ce prix, ce que tout le monde fera.
\textbf{10.} Clôtures 20.00~; 22.00~; 24.20~; 26.62~; 28.71~: rendements $+10\%$,
$+10\%$, $+10\%$, $+7.9\%$.
\textbf{11.} $\sum r^2 = 0.0362$ contre $(0.4356)^2 = 0.1898$~: cinq fois
moins.
\textbf{12.} $\sqrt{0.0362/20} = 4.3\%$ par jour, contre $\sqrt{0.1898/20} =
9.7\%$ pour le même mouvement effectué en un seul jour.
\textbf{13.} $(0.01 + 0.01 + 0.0079)/0.0362 = 0.77$~: un signal de momentum voit
la tendance la plus forte de son échantillon, sur une action que personne ne pouvait acheter.
\textbf{14.} $29.00/28.60 - 1 = 1.40\%$ en trois mois, soit 5.6~\% par an.
\textbf{15.} Moins 0.10~\% de droit et 1.00~\% de financement~: 0.30~\% en
trois mois, soit 1.2~\% par an au-dessus du taux d'intérêt. Si la finalisation est certaine,
il s'agit d'une pure prime de liquidité~; savoir si elle vaut l'immobilisation du bilan dépend
du seuil de rentabilité de la société.
\textbf{16.} $2\,000\,000/28.60 = 69\,930$~: 699 lots, 69\,900 actions.
\textbf{17.} Elle aurait ouvert près de 28.7 (dans sa monnaie) le mardi.
Pendant les trois jours bloqués, l'écart entre les deux cours mesure le
mouvement que le marché limité n'a pas encore été autorisé à faire, et non un
arbitrage~: les deux cours ne peuvent pas être livrés l'un contre l'autre.
\textbf{18.} Risque d'échec de l'opération. Avec une probabilité $p$ d'échec et une chute vers
environ 20 dans ce cas, le rendement brut de 1.40~\% doit compenser une perte espérée
de $p \times 30\%$~: à $p = 5\%$, l'opération perd déjà de l'argent.
\textbf{19.} \textbf{Trois jours.}
\textbf{20.} Elle protège les participants contre le fait de négocier à des prix formés dans la panique
ou par erreur en une journée, au prix de la suspension de la capacité de chacun à négocier
au bon prix pendant autant de jours que la nouvelle est importante.
\end{solution}

\begin{solution}{iq:m1:asian-and-emerging-equity-markets:1}
Une bande autour du cours de clôture précédent, en dehors de laquelle aucun ordre n'est accepté. Si la fair value
dépasse cette bande, l'action passe à la hausse limite ou à la baisse limite~: file à sens unique,
presque aucune négociation, et le mouvement s'achève sur plusieurs jours, la bande de chaque jour
étant recalculée à partir du cours de clôture précédent.
\iqlookfor{la file à sens unique, et le fait que le mouvement se produit malgré tout.}
\end{solution}

\begin{solution}{iq:m1:asian-and-emerging-equity-markets:2}
Volatilité (sous-estimée, puisque les grands mouvements se répartissent sur plusieurs jours) et
autocorrélation (artificiellement positive). Viennent aussi les estimations des queues de distribution et de value at risk,
les corrélations avec des marchés sans contrainte les jours d'événement, et tout backtest
qui suppose une exécution au cours de clôture d'une journée limite.
\iqlookfor{le mécanisme, le report d'un jour sur l'autre, et pas seulement la liste.}
\end{solution}

\begin{solution}{iq:m1:asian-and-emerging-equity-markets:3}
Tout ce dont la marge par aller-retour est inférieure aux 20 points de base de taxe~:
le \omterm{def:m1:what-a-trading-firm-does:market-maker}{teneur de marché} électronique, l'arbitrage sur indices et ETF, les stratégies statistiques
intrajournalières. La liquidité est alors fournie par des investisseurs à horizon plus long et par
les courtiers, les fourchettes s'élargissent, et la formation des prix migre vers les instruments
que la taxe n'atteint pas --- contrats à terme, swaps, certificats de dépôt ou \omterm{def:m1:what-a-trading-firm-does:market-maker}{teneurs de marché}
exonérés, lorsque la loi accorde des exonérations.
\iqlookfor{la migration vers les instruments non taxés.}
\end{solution}

\begin{solution}{iq:m1:asian-and-emerging-equity-markets:4}
Les deux lignes ne sont pas fongibles~: l'une ne peut pas être livrée contre l'autre,
et les investisseurs autorisés à détenir chacune diffèrent, de même que leurs coûts de financement et
les contraintes sur les ventes à découvert. Il n'existe donc pas de date de convergence. L'écart
peut être négocié comme une position statistique de valeur relative --- long de la ligne bon marché, short de la ligne chère
lorsque la vente à découvert est possible --- dimensionnée pour la possibilité que l'écart se creuse
pendant des années, ou comme un pari sur des réformes d'accès qui le réduiraient.
\iqlookfor{pas de convergence forcée, et un dimensionnement en conséquence.}
\end{solution}

\begin{solution}{iq:m1:asian-and-emerging-equity-markets:5}
Règles traitées comme des données, et non comme du code~: une table indexée par marché, classe d'instrument et
date de validité, contenant la taille de lot, la table des ticks, la formule de limite et son arrondi,
les horaires de séance et les contraintes sur les ventes~; un validateur générique unique qui la lit.
Chaque modification devient une nouvelle ligne datée, renvoyant à sa circulaire source~; les backtests
interrogent la table à la date simulée. On teste avec les exemples propres à la bourse et l'on rejoue
des ordres de production récents contre la nouvelle version de la règle avant sa date d'effet.
Les rejets émis par la bourse sont remontés
comme signal de surveillance indiquant que la table est périmée.
\iqlookfor{versionnement par date et validateur générique unique.}
\end{solution}

\begin{solution}{iq:m1:asian-and-emerging-equity-markets:6}
Le plus sévère~: le long--short actions (plus de nouvelles ventes à découvert, et les ventes existantes
peuvent être rappelées), la couverture des obligations convertibles et des bons de souscription, et le
\omterm{def:m1:what-a-trading-firm-does:market-maker}{teneur de marché} sur options sur action (on ne peut pas couvrir des puts vendus ou des calls achetés
par une vente à découvert). Viennent ensuite l'arbitrage sur indices et le \omterm{def:m1:what-a-trading-firm-does:market-maker}{teneur de marché} sur ETF
(la jambe de panier vendue), puis les revenus de prêt pour les détenteurs longs. Les contrats à terme
se négocient typiquement avec une décote par rapport à la fair value, car ils deviennent le seul moyen d'être short.
Premier matin~: vérifier si les ventes à découvert existantes peuvent être conservées~; remplacer les shorts
sur noms individuels par des contrats à terme sur indices lorsque la corrélation le permet, en acceptant
le risque de base et le risque spécifique~; réduire l'\omterm{def:m1:pnl-and-positions:exposure}{exposition brute} à ce que la couverture peut supporter~;
et intégrer dans le prix que l'interdiction n'a pas de fin annoncée.
\iqlookfor{la décote des contrats à terme comme conséquence, et les premières actions concrètes.}
\end{solution}
